\documentclass[12pt]{article}

\usepackage[left=2cm,top=2cm,right=2cm, bottom=2cm]{geometry}
\usepackage[T1]{fontenc}
\usepackage[utf8]{inputenc}
\usepackage[spanish]{babel}
\usepackage{amsfonts,setspace}
\usepackage{amsmath}
\usepackage{amssymb, amsmath, amsthm}
\usepackage{comment}
\usepackage{amssymb}
\usepackage{dsfont}
\usepackage{anysize}
\usepackage{multicol}
\usepackage{enumerate}
\usepackage{graphicx}
\usepackage{fancyhdr}
\usepackage{wasysym}
\usepackage{enumerate}
\usepackage{enumitem}
\usepackage{hyperref}
\usepackage[onelanguage, ruled, linesnumbered]{algorithm2e}

\pagestyle{fancy}

\theoremstyle{plain}
\newtheorem{teo}{Teorema}
\newtheorem{lemma}{Lemma}
\newtheorem{prop}{Proposici\'on}
\newtheorem{cor}{Corolario}
\theoremstyle{definition}
\newtheorem{defi}{Definici\'on}
\newtheorem{eje}{Ejemplo}
\newtheorem{obs}{Observaci\'on}
\newtheorem{propi}{Propiedades}

\newcommand{\norm}[1]{\lVert#1\rVert}
\newcommand{\ds}{\displaystyle}
\newcommand{\R}{\mathbb{R}}

% Error fixes
\makeatletter
\newcommand\RedeclareMathOperator{%
  \@ifstar{\def\rmo@s{m}\rmo@redeclare}{\def\rmo@s{o}\rmo@redeclare}%
}
% this is taken from \renew@command
\newcommand\rmo@redeclare[2]{%
  \begingroup \escapechar\m@ne\xdef\@gtempa{{\string#1}}\endgroup
  \expandafter\@ifundefined\@gtempa
     {\@latex@error{\noexpand#1undefined}\@ehc}%
     \relax
  \expandafter\rmo@declmathop\rmo@s{#1}{#2}}
% this is just \@declmathop without \@ifdefinable
\newcommand\rmo@declmathop[3]{%
  \DeclareRobustCommand{#2}{\qopname\newmcodes@#1{#3}}%
}
\@onlypreamble\RedeclareMathOperator
\makeatother

\DeclareMathOperator{\sen}{sen}
\RedeclareMathOperator{\cos}{cos}
\RedeclareMathOperator{\tan}{tan}
\RedeclareMathOperator{\sec}{sec}
\DeclareMathOperator{\cosec}{cosec}
\DeclareMathOperator{\cotan}{cotan}
\DeclareMathOperator{\arcsen}{arcsen}
\RedeclareMathOperator{\arccos}{arccos}
\RedeclareMathOperator{\arctan}{arctan}

\DeclareMathOperator{\senh}{senh}
\RedeclareMathOperator{\cosh}{cosh}
\RedeclareMathOperator{\tanh}{tanh}
\DeclareMathOperator{\sech}{sech}
\DeclareMathOperator{\cosech}{cosech}
\DeclareMathOperator{\cotanh}{cotanh}
\DeclareMathOperator{\arcsenh}{arcsenh}
\DeclareMathOperator{\arccosh}{arccosh}
\DeclareMathOperator{\arctanh}{arctanh}

\DeclareMathOperator{\Dom}{Dom}
\DeclareMathOperator{\Rec}{Rec}
\RedeclareMathOperator{\Im}{Im}

\DeclareMathOperator{\Int}{Int}
\DeclareMathOperator{\Adh}{Adh}
\DeclareMathOperator{\Fr}{Fr}
\DeclareMathOperator{\co}{co}

\DeclareMathOperator{\dist}{dist}
\DeclareMathOperator*{\argmin}{arg\,min}
\DeclareMathOperator*{\argmax}{arg\,max}

\let\lim=\undefined\DeclareMathOperator*{\lim}{\text{lím}}
\let\max=\undefined\DeclareMathOperator*{\max}{\text{máx}}
\let\min=\undefined\DeclareMathOperator*{\min}{\text{mín}}
\let\inf=\undefined\DeclareMathOperator*{\inf}{\text{ínf}}

\newcommand{\pint}[2]{\left< #1,#2\right>}
\newcommand{\ssi}{\Longleftrightarrow}
\newcommand{\imp}{\Longrightarrow}
\newcommand{\pmi}{\Longleftarrow}
\newcommand{\ipartial}[2]{\dfrac{\partial #1}{\partial #2}}
\newcommand{\ider}[2]{\dfrac{d #1}{d #2}}
\newcommand{\iipartial}[2]{\dfrac{\partial^2 #1}{\partial #2^2}}
\newcommand{\iider}[2]{\dfrac{d^2 #1}{d #2^2}}
\newcommand{\ijpartial}[3]{\dfrac{\partial^2 #1}{\partial #2 \partial #3}}

\newcommand{\N}{\mathbb{N}}
\newcommand{\Z}{\mathbb{Z}}
\newcommand{\Q}{\mathbb{Q}}
\newcommand{\C}{\mathbb{C}}
\newcommand{\K}{\mathbb{K}}
\renewcommand{\P}{\mathbb{P}}

\newcommand{\ALG}{\mathrm{ALG}}
\newcommand{\alg}{\mathrm{alg}}
\newcommand{\OPT}{\mathrm{OPT}}
\newcommand{\opt}{\mathrm{opt}}

\newcommand{\A}{\mathcal{A}}
\newcommand{\B}{\mathcal{B}}
\newcommand{\D}{\mathcal{D}}
\newcommand{\E}{\mathcal{E}}
\newcommand{\F}{\mathcal{F}}
\renewcommand{\H}{\mathcal{H}}
\newcommand{\I}{\mathcal{I}}
\renewcommand{\L}{\mathcal{L}}
\newcommand{\M}{\mathcal{M}}
\newcommand{\Partes}[1]{\mathcal{P}(#1)}
\renewcommand{\S}{\mathcal{S}}
\newcommand{\Tau}{\mathcal{T}}
\newcommand{\V}{\mathcal{V}}
\newcommand{\X}{\mathcal{X}}
\newcommand{\vacio}{\emptyset}
\newcommand{\tsc}{\textsc}

\newcommand{\header}{
    \fancyhead[L]{Facultad de Ciencias Físicas y Matemáticas}
	\fancyhead[R]{Universidad de Chile}
	\vspace*{-1cm}
	\hspace{-0.5cm}\begin{minipage}{0.7\textwidth}
	\begin{flushleft}
	\textbf{MA3705 Algoritmos Combinatoriales}.\\
	\textbf{Profesor:} Iván Rapaport.\\
	\textbf{Auxiliares:} Antonia Labarca y Cristian Palma.\\
	\end{flushleft}
	\end{minipage}
	\begin{minipage}{0.3\textwidth}
	\begin{flushright}
	\hspace{-0.5cm}\includegraphics[scale=0.2]{img/dim.pdf}
	\end{flushright}
	\end{minipage}

	\medskip
}
\begin{document}
\header
\begin{center}
	\LARGE \bf{Auxiliar 5}
\end{center}

Suponga que cuenta con un procesador \textit{monocore} (es decir, que solo puede realizar una tarea a la vez).
Además, tiene un conjunto finito $T = \{a_1, a_2, \dots, a_n\}$ de posibles tareas a realizar, todas con duración de 1 unidad de tiempo.
Cada una de las tareas tiene asociado un plazo $d_i$ tal que $\forall i \in [n], 1\leq d_i \leq n$, tal que si la tarea $a_i$ se concluye a lo más en el tiempo $d_i$ se recibe una ganancia $w_i$.

Diremos que un \textit{schedule} $S$ de $T$ es un ordenamiento de todas las tareas de $T$ (note que hay $n!$ \textit{schedules} posibles). Se considerará que la primera tarea inicia en el tiempo 0 y termina en el tiempo 1, la segunda inicia en el tiempo 1 y termina en el tiempo 2, \dots, la $n$-ésima inicia en el tiempo $n-1$ y termina en el tiempo $n$.

Decimos que para un determinado \textit{schedule} $S$, una tarea $a_i$ está \textbf{a tiempo} si está en la posición a lo más $d_i$ del ordenamiento (por lo que se completaría a lo más en el tiempo $d_i$). Por el contrario, está \textbf{atrasada} si está en una posición posterior.

Sean $E(S)$ el conjunto de las tareas a tiempo y $L(S)$ el conjunto de las tareas atrasadas.
Diremos que un \textit{schedule} está en \textbf{forma normal} si para algún $0\leq j\leq n$ las primeras $j$ tareas están a tiempo y las siguientes $n-j$ están atrasadas. Diremos que está en \textbf{forma canónica} si está en forma normal y además las tareas a tiempo están ordenadas ascendentemente según su deadline.


\begin{enumerate}[label ={\bf P\arabic*}]
	\item Justifique que para $S$ se puede construir $\bar{S}$ en forma canónica de forma que $E(S)=E(\bar{S}), L(S)=L(\bar{S})$.
\end{enumerate}
    
    Decimos que un conjunto $I$ de tareas es independiente si existe un \textit{schedule} $S$ tal que $I\subseteq E(S)$. Sea $\mathcal{I} =\{A \subseteq T | A \text{ es independiente}\}$.\\

    Para $0\leq t \leq n, A \subseteq T$, sea $N_t(A) = |\{a_i \in A | d_i \leq t\}|$.

\begin{enumerate}[label={\bf P\arabic*}]\addtocounter{enumi}{1}
    \item Pruebe que para $A\subseteq T$ son equivalentes \begin{enumerate}
        \item $A\in\mathcal{I}$
        \item $\forall t=0\dots n, N_t(A)\leq t$
        \item Si se ordena $a$ ascendentemente según el plazo de cada tarea, ninguna tarea estará atrasada.
    \end{enumerate}

    \item Se define el sistema de independencia $(T, \mathcal{I})$. Pruebe que es una matroide.
    
    \item Diseñe y analice un algoritmo \textit{greedy} que resuelva el problema.

\end{enumerate}
\end{document}